\section{Production Successor: Complete Analysis in One Hop}
\label{sec:one-hop}
The production \code{SpectrumAnalysis<T>} replaces the modules' original
completion-driven loop with one ordered sequence per frame. It computes the
same windowed transform and magnitude processing while omitting unused
conjugate bins. The Fourier module instantiates four independent SIMD channels;
the Spectre module uses scalar samples. The current implementation also
weights window-cache rebuilding and Fourier's coordinate output as described
below. The historical timing
campaigns retained in this report predate that integration; their measurements
must not be attributed to the revised scheduler.

\subsection{Dependency order and work bound}
Let $M=N/2$, $K=M+1$, and $B=(M/2)\log_2 M$. Table~\ref{tab:units} defines
the operations in each stage. Let $p=4$ while rebuilding window coefficients
and $p=1$ for a clean window cache. Let $o$ be the compile-time output weight:
one by default, or two for Fourier's four-channel coordinate mapping. The
ordered scheduling sequence has length
\begin{equation}
 W=pM+B+(1+o)K=\frac{N}{4}\log_2(N/2)+\frac{(p+1+o)N}{2}+1+o.
 \label{eq:wholework}
\end{equation}
All packing completes before any butterfly, all butterflies before any
reconstruction, and the complete prefix sum before any range smoothing. A
sample quota may cross stage boundaries. The dependency graph is unchanged
when execution pauses; it therefore admits the same induction argument as
the butterfly traversal in Appendix~\ref{sec:legacy-transform}. This is an argument about preserving operations for
a chosen implementation, not bitwise identity to older code with different
twiddle precision or compiler transformations.

\begin{table}[tb]
\centering\small
\begin{tabular}{@{}rlp{0.64\linewidth}@{}}
\toprule
Count & Stage & Work performed by one unit \\
\midrule
$pM$ & Pack & Each pair occupies $p$ units: read two retained samples, prepare/apply window weights and store the permuted pair at the first unit; skip the remaining credit. \\
$B$ & Transform & Execute one complex butterfly and advance traversal indices. \\
$K$ & Reconstruct & Recover one bin, take its magnitude, and append a prefix sum if frequency smoothing is enabled. \\
$oK$ & Output & Each bin occupies $o$ units: apply enabled smoothing and invoke the output callback at the first unit; skip the remaining credit. \\
\bottomrule
\end{tabular}
\caption{Production scheduling units. Dirty window preparation has weight four; Fourier's coordinate output has weight two. Other operations have weight one. These are bookkeeping units, not equal-duration operations.}
\label{tab:units}
\end{table}


For call $s\in\{1,\ldots,H\}$, assign the quota
\begin{equation}
 q_s=\floor{sW/H}-\floor{(s-1)W/H}.
 \label{eq:balanced}
\end{equation}
The cumulative work after $s$ calls is $\floor{sW/H}$, so completion is
exactly on call $H$, with at most $\ceil{W/H}$ units per call. For
$H=\rho N$ with fixed $\rho>0$, this is $O(\log N)$. A quotient/remainder
accumulator implements the rule without per-sample division. This elementary
balanced schedule is not a new scheduling result; its use here includes the
complete dependency-ordered analysis. Appendix~\ref{app:one-hop} preserves
the dispatch listing, proof, and clean/dirty-cache examples.

The bound includes scheduled window/cache rebuilds and bounded module output
callbacks. Construction and host lifecycle callbacks are excluded; constant
input/control work and an $O(\log N)$ plan-index calculation remain. No
$O(N)$ boundary pass remains in steady-state engine analysis. Reweighting
redistributes work but does not make units equal in duration, constrain an
arbitrary user callback, or bound operating-system interruptions.

\subsection{Retained samples, plans, and publication time}
Let the frame endpoint be $t_r=rH$, counting continuously captured samples
from zero. Frame selection occurs on call $t_r$, and the last scheduling
credit and publication occur on call $t_r+H-1$. The last actual output may
precede its final credit; it remains private until publication. Consequently
\begin{equation}
 A_{\mathrm{newest}}=(H-1)/f_s,\qquad
 A_{\mathrm{mid}}=((N-1)/2+H-1)/f_s.
 \label{eq:wholeage}
\end{equation}
Successive publications are exactly $H$ samples apart. This convention differs
from the original next-call publication in Equation~\eqref{eq:age}. Startup
uses zero padding; a full first window is not awaited. With changing hops the
endpoints satisfy $t_{r+1}=t_r+H_r$, and each frame uses its latched $H_r$.

Incremental packing must retain the selected samples while capture continues.
A preallocated ring of capacity $C=N_{\max}+H_{\max}$ suffices for the
prepared maximum window and hop; no whole-frame capture copy is needed.
The lifetime proof and maximum-size table reuse are retained in
Appendices~\ref{app:one-hop} and~\ref{app:implementation}.

\subsection{Smoothing, configuration, and output ownership}
Reconstruction obtains magnitudes and, when frequency smoothing is enabled,
builds a prefix sum before any range average is emitted. Output applies an
inclusive fractional-octave magnitude mean, optional temporal averaging, and
the module's display conversion. These are magnitude means, not power means
or constant-Q analysis. The window normalization and display units are
preserved. Appendix~\ref{app:one-hop} records the equations, control
conventions, cache behavior, and sample-level hop quantization.

Settings latch at each frame start; reset and sample-rate changes cancel
partial analysis. Window and smoothing caches rebuild inside scheduled units,
and length changes logically clear input and averaging history. Storage is
prepared outside the sample path; increasing the supported hop in a
sample-rate callback can allocate. Freeze behavior differs between the two
modules, so continuous-input timestamp formulas do not apply across pauses.

Each output bin is written to a producer-owned destination. Only a complete
frame is published through a single-producer/single-consumer triple-slot
mailbox using acquire/release atomic exchange. The consumer retains its slot
until its next successful exchange, and the producer never overwrites that
slot. A slow consumer may miss intermediate snapshots. Fourier uses one curve
mailbox; Spectre uses one mailbox per history column, with image history and
recoloring on the UI thread. Lifecycle work, other shared controls, UI work,
and host scheduling remain outside the quota; the contract is not a broad
claim of hard real-time safety.
